\section*{Capítulo \ref{ch:b1:p-block-16-18} --- O bloco p II: oxigênio, halogênios e gases nobres}

\begin{solution}{exo:b1:p-block-16-18:1}
\ce{H2S} $-$II; \ce{S8} 0; \ce{SO2} $+$IV; \ce{SO3^2-} $+$IV; \ce{SO4^2-}
$+$VI; \ce{S2O3^2-} $+$II em média; \ce{H2S2O7} $+$VI.
\end{solution}

\begin{solution}{exo:b1:p-block-16-18:2}
Um \omterm{def:b1:p-block-16-18:halogen}{halogênio} oxida os haletos situados abaixo dele: \ce{Cl2 + 2I- -> 2Cl- + I2} e
\ce{Cl2 + 2Br- -> 2Cl- + Br2} ocorrem; o iodo com o brometo e o bromo com
o cloreto, não.
\end{solution}

\begin{solution}{exo:b1:p-block-16-18:3}
\ce{SO2}: S com um \omterm{def:b1:lewis-resonance-vsepr:lewis}{par não ligante}, uma S=O e uma \ce{S+-O-} em cada uma
de duas \omterm{def:b1:lewis-resonance-vsepr:resonance}{estruturas de ressonância} (ou duas S=O com octeto expandido):
angular. \ce{SO3}: trigonal plana, três ligações S--O equivalentes.
\ce{O3}: O central com um \omterm{def:b1:lewis-resonance-vsepr:lewis}{par não ligante}, uma O=O e uma O--O$^-$, duas
\omterm{def:b1:lewis-resonance-vsepr:resonance}{estruturas de ressonância}: angular.
\end{solution}

\begin{solution}{exo:b1:p-block-16-18:4}
\ce{SO2 + H2O <=> H2SO3}; \ce{SO2 + 2NaOH -> Na2SO3 + H2O};
\ce{SO3 + H2O -> H2SO4}; \ce{SO3 + 2NaOH -> Na2SO4 + H2O}.
\end{solution}

\begin{solution}{exo:b1:p-block-16-18:5}
HF: $x^2/(0.10 - x) = 10^{-3.16}$, $x = \qty{8.0e-3}{mol/L}$, pH 2.10. HCl:
pH 1.00. A ligação H--F é forte e o íon fluoreto forma \omterm{def:b1:intermolecular-forces-solvents:hbond}{ligações de
hidrogênio} fortes: HF é um \omterm{def:b1:predominant-reaction:strong-acid}{ácido fraco}.
% ledger: dfg:HF_ao, dfg:F-_ao
\end{solution}

\begin{solution}{exo:b1:p-block-16-18:6}
\ce{BrF5}: cinco ligações e um \omterm{def:b1:lewis-resonance-vsepr:lewis}{par não ligante}, pirâmide de base
quadrada. \ce{IF7}: bipirâmide pentagonal. \ce{ICl4-}: quatro ligações e
dois \omterm{def:b1:lewis-resonance-vsepr:lewis}{pares não ligantes}, quadrado-plana. \ce{XeO3}: três ligações e um \omterm{def:b1:lewis-resonance-vsepr:lewis}{par
não ligante}, pirâmide trigonal.
\end{solution}

\begin{solution}{exo:b1:p-block-16-18:7}
$\log K = 2(1.36 - 0.53)/0.059 = 28$. O iodo liberado forma com o amido o
\omterm{def:b1:complexation:complex}{complexo} azul intenso.
\end{solution}

\begin{solution}{exo:b1:p-block-16-18:8}
$E^\circ(\ce{O3}/\ce{O2}) = 2.07 > 0.53$~V: \ce{O3 + 2I- + 2H+ -> O2 + I2
+ H2O}. O iodo formado é titulado com tiossulfato
(\cref{exo:b1:nernst:7}): um \ce{I2} por \ce{O3}.
\end{solution}

\begin{solution}{exo:b1:p-block-16-18:9}
\ce{C12H22O11 -> 12C + 11H2O}; $M = \qty{342.0}{g/mol}$, $10.0/342.0 =
\qty{0.0292}{mol}$, carbono $0.0292 \times 12 \times 12.0 = \qty{4.21}{g}$.
\end{solution}

\begin{solution}{exo:b1:p-block-16-18:10}
$x(0.10 + x)/(0.10 - x) = 10^{-1.99}$: $x = \qty{8.6e-3}{mol/L}$,
$[\ce{H3O+}] = 0.1086$, $\mathrm{pH} = 0.96$ em vez de 1.00: a segunda
acidez acrescenta cerca de \qty{9}{\percent} a $[\ce{H3O+}]$.
\end{solution}

\begin{solution}{exo:b1:p-block-16-18:11}
\ce{F2}/\ce{F-} (2.89~V) fica acima de todos os outros pares: nenhum
\omterm{def:b1:nernst:couple}{oxidante} pode retirar o elétron do fluoreto. Na eletrólise aquosa, a
água, oxidada a um potencial muito mais baixo (\ce{O2}/\ce{H2O}, 1.23~V),
reagiria em seu lugar; um fluoreto fundido não contém água.
\end{solution}

\begin{solution}{exo:b1:p-block-16-18:12}
$\Delta_r G^\circ = -51.4 - (16.40 - 51.57) = \qty{-16.2}{kJ/mol}$,
$\log K = 16.2/5.708 = 2.84$, $K = 7 \times 10^2$.
% ledger: dfg:I2_ao, dfg:I-_ao, dfg:I3-_ao
\end{solution}

\begin{solution}{pb:b1:p-block-16-18:1}
\textbf{1.} 0, $+$IV, $+$VI, $+$VI.
\textbf{2.} \ce{S + O2 -> SO2}.
\textbf{3.} S com dois domínios ligantes e um \omterm{def:b1:lewis-resonance-vsepr:lewis}{par não ligante}: angular,
cerca de \ang{120}.
\textbf{4.} $10^6/32.1 = \qty{3.12e4}{mol}$.
\textbf{5.} \qty{3.115e4}{mol} de \ce{O2}: $V = 3.115 \times 10^4 \times
8.314 \times 298.15/10^5 = \qty{772}{m^3}$, ar $772/0.21 = \qty{3.7e3}{m^3}$.
\textbf{6.} A combustão é muito exotérmica; o \omterm{def:b1:elementary-steps:catalyst}{catalisador} funciona em uma
faixa de temperatura limitada, e a oxidação do \ce{SO2}, também
exotérmica, é menos completa a quente.
\textbf{7.} \ce{2SO2 + O2 <=> 2SO3}.
\textbf{8.} Uma constante de equilíbrio grande nada diz sobre a
velocidade: à temperatura ambiente, a reação não avança.
\textbf{9.} Ele oferece um caminho de \omterm{def:b1:rate-laws:activation-energy}{energia de ativação} mais baixa e é
regenerado: um \omterm{def:b1:elementary-steps:catalyst}{catalisador} heterogêneo.
\textbf{10.} Trigonal plana.
\textbf{11.} Mais oxigênio mantém o quociente
$Q = p_{\ce{SO3}}^2/(p_{\ce{SO2}}^2 p_{\ce{O2}})$ mais baixo: a conversão
do \ce{SO2} vai mais longe.
\textbf{12.} \qty{0.3}{\percent}.
\textbf{13.} \ce{SO3 + H2SO4 -> H2S2O7}.
\textbf{14.} Ele forma uma névoa de gotículas finas de ácido que atravessa
o absorvedor.
\textbf{15.} \ce{H2S2O7 + H2O -> 2H2SO4}.
\textbf{16.} \ce{2S + 3O2 + 2H2O -> 2H2SO4}.
\textbf{17.} Uma água por enxofre: $\num{3.115e4} \times 18.0 =
\qty{561}{kg}$.
\textbf{18.} A diluição libera muito calor; a água despejada sobre o
ácido ferve de imediato e projeta ácido.
\textbf{19.} $x(0.050 + x)/(0.050 - x) = 10^{-1.99}$: $x = \qty{7.5e-3}{mol/L}$,
$[\ce{H3O+}] = \qty{0.0575}{mol/L}$, $\mathrm{pH} = 1.24$.
\textbf{20.} $3.12 \times 10^4 \times 0.995 \times 98.1 = \qty{3.04}{t}$.
\textbf{21.} $84 \times 98.1/32.1 = \qty{257}{Mt}$.
\textbf{22.} $98.1/32.1 =$ \textbf{\qty{3.06}{t} de ácido sulfúrico por
tonelada de enxofre}.
\end{solution}
